
%\section*{PARTIE 3. ETUDE GEOMETRIQUE ET CINEMATIQUE}
\section{Etude géométrique et cinématique\label{2026_ats_s3}}

\begin{obj}
Pré-dimensionner en vitesse le moteur de l'unité de fermeture du moule et vérifier une donnée constructeur.

\textbf{Exigence 1211} : le moule a une longueur d'au moins \SI{400}{mm}.
\end{obj}
%L'unité de fermeture du moule est constituée d'un mécanisme à genouillère dont le schéma cinématique complet ainsi que le fonctionnement sont décrits dans les annexes~\ref{2026_ats_ann2} et \ref{2026_ats_ann3}. 
L'unité de fermeture du moule est constituée d'un mécanisme à genouillère  dont le schéma cinématique plan retenu pour l'étude géométrique est le suivant (figure~\ref{2026_ats_fig7}).

\begin{figure}[h]
\centering
  \includegraphics[width=\textwidth]{fig_07}
%\captionsetup{labelformat=empty}
\caption{Schéma cinématique et paramétrage \label{2026_ats_fig7}}%Figure 7 : 
\end{figure}

On donne :

\begin{minipage}[c]{.55\linewidth}
\begin{itemize}
  \item $\overrightarrow{O A}=\lambda \vec{x}_{0}+h  \vec{y}_{0} \quad$ avec $\lambda$ un paramètre algébrique (peut être $\geq 0$ ou $\leq 0$);
  \item $\overrightarrow{O B}=e  \vec{x}_{3}$;
  \item $\overrightarrow{A B}=d  \vec{x}_{2}$;
\end{itemize}
\end{minipage}
\hfill
\begin{minipage}[c]{.4\linewidth}
\begin{itemize}
  \item $\overrightarrow{O C}=2 e \vec{x}_{3}$;
  \item $\beta=\left(\vec{x}_{0}, \vec{x}_{2}\right)=\left(\vec{y}_{0}, \vec{y}_{2}\right)$;
  \item $\theta=\left(\vec{x}_{0}, \vec{x}_{3}\right)=\left(\vec{y}_{0}, \vec{y}_{3}\right)$;
  \item $\varphi=\left(\vec{x}_{0}, \vec{x}_{4}\right)=\left(\vec{y}_{0}, \vec{y}_{4}\right)$.
\end{itemize}
\end{minipage}




\section{Etude géométrique générale \label{2026_ats_s3_1}}
%Question 9. 
\begin{question}
Réaliser les figures de changements de bases mettant en évidence les paramètres angulaires.
\end{question}
\ifprof
\begin{corrige}
\end{corrige}
\else
\fi



\begin{question}
Etablir une relation scalaire ne contenant indépendante de l'angle $\beta$ et montrer qu'elle peut se mettre sous la forme suivante : 
$P  \cos \theta+Q  \sin \theta=R$ en précisant les expressions de $P$, $Q$ et $R$.
\end{question}
\ifprof
\begin{corrige}
\end{corrige}
\else
\fi
%%Question 10. 
%\begin{question}
%Ecrire une équation vectorielle de fermeture géométrique sur le cycle $OABO$. En projetant sur $\overrightarrow{x}_{0}$ et $\overrightarrow{y}_{0}$, en déduire les deux équations scalaires liant les paramètres $\lambda$, $\theta$ et $\beta$.
%\end{question}
%\ifprof
%\begin{corrige}
%\end{corrige}
%\else
%\fi
%
%%Question 11. 
%\begin{question}
%A partir des deux équations précédemment obtenues, établir une relation scalaire ne contenant plus l'angle $\beta$ et montrer qu'elle peut se mettre sous la forme suivante : 
%$P  \cos \theta+Q  \sin \theta=R$ en précisant les expressions de $P$, $Q$ et $R$.
%\end{question}
%\ifprof
%\begin{corrige}
%\end{corrige}
%\else
%\fi

Le terme $P^{2}+Q^{2}$ ne s'annulant jamais, la relation précédente peut aussi s'écrire :
$
\dfrac{P}{\sqrt{P^{2}+Q^{2}}}  \cos \theta+\dfrac{Q}{\sqrt{P^{2}+Q^{2}}} \sin \theta=\dfrac{R}{\sqrt{P^{2}+Q^{2}}}
$.

On pose : $\cos \psi=\frac{P}{\sqrt{P^{2}+Q^{2}}}$ et $\sin \psi=\frac{Q}{\sqrt{P^{2}+Q^{2}}}$. On rappelle que, pour $x$ et $y$ réels : $\cos \left(x-y\right)=\cos x  \cos y+\sin x \sin y$.

%Question 12. 
\begin{question}
Exprimer $\theta$ en fonction de $P$, $Q$ et $R$.
\end{question}
\ifprof
\begin{corrige}
\end{corrige}
\else
\fi

%3.2. 
\subsection{Prédimensionnement en vitesse du moteur \label{2026_ats_s3_2}}
L'expression trouvée précédemment étant trop complexe à exploiter, on fournit un tracé de la fonction $\theta=\mathrm{f}(\lambda)$ sur la figure \ref{2026_ats_fig8}.

\begin{figure}[h]
\centering
  \includegraphics[width=\textwidth]{fig_08}
%\captionsetup{labelformat=empty}
\caption{Tracé de la fonction $\theta=f(\lambda)$ \label{2026_ats_fig8}}
%Figure 8 : 
\end{figure}


Lorsque le moule est en position fermée, l'angle $\theta$ vaut $\theta=0^{\circ}$.\\
Lorsque le moule est en position ouverte, l'angle $\theta$ vaut $\theta=90^{\circ}$.

%Question 13. 
\begin{question}
Déterminer $\Delta \lambda$, l'amplitude du déplacement horizontal de la pièce 1 lors d'une ouverture (ou fermeture) de moule.
\end{question}
\ifprof
\begin{corrige}
\end{corrige}
\else
\fi

%Question 14. 
\begin{question}
Exprimer puis calculer en \si{m/s} la vitesse moyenne $\dot{\lambda}_{\text {moy }}$ de déplacement de la pièce 1 si on souhaite que ce déplacement soit effectué en un temps $\Delta t=\SI{1,2}{s}$ (on néglige ici les phases d'accélération et de décélération).
\end{question}
\ifprof
\begin{corrige}
\end{corrige}
\else
\fi


Le dispositif d'entraînement en translation de l'écrou 1 est constitué d'un moteur électrique, d'une transmission poulies/courroie crantées et d'un système vis/écrou $\underline{\mathbf{6}} / \mathbf{1}$ dont le pas vaut $p_{v}=\SI{0,03}{m/tr}$ (voir figure \ref{2026_ats_fig9}).

On note $\mathrm{N}_{\mathrm{m}}$ la vitesse de rotation du moteur et $\mathrm{N}_{6 / 0}$ la vitesse de rotation de la vis $\underline{6}$ par rapport au bâti 0. Ces vitesses sont exprimées en $\mathrm{tr} / \mathrm{min}$.

\begin{figure}[h]
\centering
  \includegraphics[width=.8\textwidth]{076c1ba0-5a37-41d5-876a-83fc91dd5c32-10_1015_1549_1131_284}
%\captionsetup{labelformat=empty}
\caption{Schéma de la transmission de puissance \label{2026_ats_fig9}}%Figure 9 : 
\end{figure}

%Question 15. 
\begin{question}Etablir la relation entre la vitesse moyenne $\dot{\lambda}_{\text {moy }}$ de déplacement de la pièce 1 et la vitesse de rotation $\mathrm{N}_{6 / 0}$ de la vis $\underline{6}$ en faisant intervenir le pas $\mathrm{p}_{\mathrm{v}}$ du système vis/écrou 6/1.
\end{question}
\ifprof
\begin{corrige}
\end{corrige}
\else
\fi

%Question 16. 
\begin{question}
Exprimer la vitesse de rotation $N_{m}$ du moteur en fonction de $N_{6 / 0}, Z_{6}$ et $Z_{m}$. En déduire $\mathrm{N}_{\mathrm{m}}$ en fonction de $\dot{\lambda}_{\text {moy }}, \mathrm{p}_{\mathrm{v}}, \mathrm{Z}_{6}$ et $\mathrm{Z}_{\mathrm{m}}$. Calculer la valeur numérique de $\mathrm{N}_{\mathrm{m}}$ (en $\mathrm{tr} / \mathrm{min}$) permettant d'obtenir l'ouverture du moule dans le temps souhaité.
\end{question}
\ifprof
\begin{corrige}
\end{corrige}
\else
\fi

%3.3. 
\subsection{Détermination de l'amplitude d'ouverture du moule \label{2026_ats_s3_3}}
Une grande ouverture du moule est intéressante pour la fabrication car cela libère de l'espace pour l'éjection de pièces ayant une profondeur élevée.

L'angle $\theta$ précédemment exprimé intervient dans le système de transformation de mouvement de type bielle-manivelle constitué des pièces $\underline{\mathbf{3}}$, $\underline{\mathbf{4}}$ et $\underline{\mathbf{5}}$. Les positions extrêmes du moule $\mathbf{5}$ sont obtenues pour $\theta=0^{\circ}$ et $\theta=90^{\circ}$.

On donne les valeurs numériques suivantes : $e = \SI{150}{mm}$ et $L=\SI{400}{mm}$.


%Question 17. 
\begin{question}
A partir de l'étude du triangle OCD (figure \ref{2026_ats_fig7}) pour $\theta=90^{\circ}$, déterminer l'expression puis la valeur numérique de $\|\overrightarrow{O D}\|$.
\end{question}
\ifprof
\begin{corrige}
\end{corrige}
\else
\fi

%Question 18.
\begin{question}
En menant une étude similaire, déterminer l'expression puis la valeur numérique de $\|\overrightarrow{O D}\|$ pour $\theta=0^{\circ}$. En déduire l'amplitude $\Delta x$ d'ouverture du moule $\underline{\mathbf{5}}$. Conclure quant au respect de l'exigence '1211'.
\end{question}
\ifprof
\begin{corrige}
\end{corrige}
\else
\fi
%
%\begin{figure}[h]
%\centering
%  \includegraphics[width=\textwidth]{fig_ann_02.png}
%%\captionsetup{labelformat=empty}
%\caption{Annexe 2 : Schéma cinématique complet \label{2026_ats_ann2}}%Figure 9 : 
%\end{figure}
%
%
%\begin{figure}[h]
%\centering
%  \includegraphics[width=\textwidth]{fig_ann_03.png}
%%\captionsetup{labelformat=empty}
%\caption{Annexe 3 : Phases de fonctionnement\label{2026_ats_ann3}}%Figure 9 : 
%\end{figure}
